\documentclass[12pt]{article}
\usepackage[utf8]{inputenc}
\usepackage[T1]{fontenc}
\usepackage{hyperref}
\title{Ugaritic Language Resources for K-12 Curriculum Development}
\author{Compiled by Claude Sonnet 5 (claude-sonnet-5) for Andr\'as Kornai}
\date{}
\begin{document}
\maketitle
\section{Introduction}
Ugaritic is an extinct Northwest Semitic language (Afro-Asiatic) known from alphabetic cuneiform tablets from Ugarit. It has no native speakers, so any curriculum would be heritage and reading-based, not conversational. The search turned up mostly library catalogue records. It found no free online corpus and no school materials.
\section{Primary Educational Resources}
No K-12 primers or school textbooks were found. Simon Parker's \emph{Ugaritic Narrative Poetry} (SBL, 1997) gives transcriptions and English translations of many texts and could serve as teacher background. The search results did not include a URL for it.
\section{Linguistic and Reference Materials}
\begin{itemize}
\item Bordreuil and Pardee, \emph{A Manual of Ugaritic}: a student-oriented grammar. The English edition includes a CD with a hyperlinked PDF. Catalogue record: \url{https://oeaw.ac.at/resources/Record/9781575066523}
\item Catalogue record for \emph{A Manual of Ugaritic}: \url{https://search.ub.tum.de/vufind/Record/DE-604.BV047416306}
\item Lexicity wiki page on Ugaritic, listing some free aids such as a verb forms chart. Its dictionary and text sections were empty: \url{https://oderalon.net/wiki/index.php/Lexicity/Ugaritic}
\end{itemize}
Sivan's \emph{A Grammar of the Ugaritic Language} (Brill, 1997), Tropper's \emph{Ugaritische Grammatik}, and the del Olmo Lete and Sanmartin dictionary (Brill) were named in search results without usable URLs. The standard text edition, Dietrich, Loretz and Sanmartin's \emph{The Cuneiform Alphabetic Texts} (3rd ed., 2013), was also named without a free online version.
\section{Religious and Community Resources}
None found in the search results.
\section{Digital and Audio Resources}
No audio or free digital corpus was found. Searches for KTU/CAT databases were not completed.
\section{Curriculum Development Context}
Ugaritic is documented through scholarly grammars, dictionaries and text editions aimed at university students. A K-12 programme would need to adapt these materials, for example into a script-learning unit or a unit on translated myths. No existing school programme was found.
\end{document}